\FloatBarrier
\section{Supplementary comparison with a single classifier}
\label{sec:singlebaseline}
\subsection{Question and comparison protocol}
Could one classifier reproduce the hybrid's balance of family detection and background errors? This supplementary experiment was specified during manuscript review, after the hybrid results were known. It reuses the final training, calibration and evaluation collections and keeps the selected hybrid models. It tests the need for the compound classifier and alarm path within the retained representation \cite{projectbaseline}.

Each control uses one LightGBM binary classifier trained jointly on the five attack families. The normal reference is retained: ``single classifier'' refers to the detection stage, not to an unfitted preprocessing pipeline. Inputs are the 116 causal features in Modena and those features plus the 42 received-pressure-history columns in L-Town. The zero-hour control uses only these causal inputs. The three-hour control adds the delayed heads' 59 bounded-window features, giving 175 and 217 columns, respectively. Temporal features are computed from complete received series before training endpoints are sampled. No expert predictions, routing outputs or family labels enter either control at inference.

The main complexity comparison is the three-hour control, which has the same maximum information allowance as the hybrid. It makes one delayed decision per endpoint using one score and threshold. The hybrid also keeps General's immediate path. Both receive the same observation types and have the same maximum deadline, but the hybrid uses learned intermediate scores and aggregated evidence. The comparison measures the complete designs under that allowance, rather than the router alone.

All ten model seeds are fitted for both controls in both networks. Each fit retains every positive TRAIN endpoint and samples up to 60,000 negative endpoints using the existing General sampling convention. The loss gives half its weight to positive endpoints, equally across families and events, and half to negatives, equally across sampled scenarios. The two controls share the sampled rows and weights within a seed. The fixed LightGBM recipe uses 400 trees, 31 leaves, learning rate 0.03, minimum leaf size 20, L1/L2 penalties of 1 and 20, and feature fraction 0.85. These settings follow the existing full temporal classifier recipe; no hyperparameter search is added.

Calibration searches distinct score thresholds under the existing clean-FPR ceiling of 0.005. It must return two operating points: one maximizing pooled F1 and one maximizing worst-family F1. Ties are broken first by the other F1 criterion, then by lower clean FPR. Pooled F1 defines the primary comparison; the family-balanced objective tests the threshold trade-off. These explicit single-threshold alternatives differ from the hybrid's recorded integration and protection searches.

All 40 classifiers and 80 operating points are frozen before supplementary evaluation on the same six source datasets used for each network's final hybrid result. Means and paired differences use the 60 model/source cells defined in Equations~\ref{eq:cellmean} and~\ref{eq:paireddifference}. Both time allowances and calibration objectives are reported, and the locked test remains unused. Because these sources have already been assessed, this is a supplementary controlled comparison rather than a fresh independent confirmation.

\subsection{Performance and complexity trade-offs}
\begin{table}[!htbp]
\centering\small
\caption[Single-classifier comparison on Modena]{Modena: single-classifier controls versus the frozen hybrid, each averaged over ten model seeds and six shared sources. Single-classifier thresholds maximize calibration pooled F1 under the common clean-FPR ceiling; the hybrid retains its recorded operating points. Clean FPR is a percentage.}
\label{tab:singlemodena}
\begin{tabular}{@{}lrrr@{}}
\toprule
Metric & Single, 0 h & Single, 3 h & Hybrid\\
\midrule
Random F1 & 0.9672 & 0.9740 & 0.9549\\
Replay F1 & 0.9278 & 0.9060 & 0.9283\\
Drift F1 & 0.7015 & 0.7627 & 0.8274\\
Noise F1 & 0.6135 & 0.6742 & 0.8439\\
Targeted F1 & 0.9731 & 0.9795 & 0.9680\\
Family macro F1 & 0.8366 & 0.8593 & 0.9045\\
Pooled F1 & 0.8259 & 0.8509 & 0.8669\\
Clean FPR (\%) & 0.0374 & 0.0317 & 0.0584\\
\bottomrule
\end{tabular}
\end{table}
\FloatBarrier
In Modena, the three-hour single classifier reaches pooled F1 of 0.8509 versus 0.8669 for the hybrid. Its drift and noise means are 0.7627 and 0.6742, compared with 0.8274 and 0.8439 for the hybrid. Replay rises from 0.9060 to 0.9283. Family macro F1 increases from 0.8593 to 0.9045, and every hybrid family mean stays above 0.80. The single classifier has higher random and targeted means and lower clean FPR: 0.0317\% versus 0.0584\%. At these operating points, the hybrid gains temporal-family coverage and pooled performance at the cost of more clean false alarms.

Selecting the family-balanced threshold for the three-hour Modena classifier raises drift and noise to 0.8311 and 0.8141. Replay falls to 0.7772 and pooled F1 to 0.7787, with clean FPR rising to 0.1921\%. The two prespecified objectives trade different parts of the family profile against each other. The hybrid keeps replay and pooled performance together with strong drift and noise means. Table~\ref{tab:balancedbaseline} reports the full alternative-objective results.

\begin{table}[!htbp]
\centering\small
\caption[Single-classifier comparison on L-Town]{L-Town: single-classifier controls versus the frozen hybrid, each averaged over ten model seeds and six shared sources. Single-classifier thresholds maximize calibration pooled F1 under the common clean-FPR ceiling; the hybrid retains its recorded operating points. Clean FPR is a percentage.}
\label{tab:singleltown}
\begin{tabular}{@{}lrrr@{}}
\toprule
Metric & Single, 0 h & Single, 3 h & Hybrid\\
\midrule
Random F1 & 0.9927 & 0.9950 & 0.9920\\
Replay F1 & 0.6927 & 0.7084 & 0.7312\\
Drift F1 & 0.7949 & 0.8698 & 0.8990\\
Noise F1 & 0.7706 & 0.8422 & 0.9279\\
Targeted F1 & 0.9935 & 0.9961 & 0.9918\\
Family macro F1 & 0.8489 & 0.8823 & 0.9083\\
Pooled F1 & 0.8685 & 0.8977 & 0.8950\\
Clean FPR (\%) & 0.0192 & 0.0184 & 0.0550\\
\bottomrule
\end{tabular}
\end{table}
\FloatBarrier
In L-Town, the three-hour single classifier's pooled F1 is 0.8977, slightly above the hybrid's 0.8950. Clean FPR is also lower, at 0.0184\% versus 0.0550\%. The hybrid has higher replay, drift and noise means: 0.7312 versus 0.7084, 0.8990 versus 0.8698, and 0.9279 versus 0.8422, respectively. These gains raise family macro F1 from 0.8823 to 0.9083. Random and targeted favour the single classifier by 0.0030 and 0.0043. The pooled means are close, but the hybrid covers the weaker families better and produces more background alarms.

Family-balanced calibration raises the L-Town control's drift and noise means to 0.8800 and 0.8653. Replay is 0.7085, pooled F1 is 0.8929 and clean FPR is 0.0325\%. Some gaps narrow, but replay, drift and noise remain below the hybrid. Under either frozen objective, the choice is between competitive pooled F1 with fewer clean alarms and better coverage of the difficult mechanisms.

\figfile{single_model_paired.pdf}{1.0}{Paired hybrid effects relative to the three-hour single classifier}{Family-F1 differences between the frozen hybrid and the three-hour single classifier at its pooled-F1 operating point. Each circle averages ten model seeds within one evaluation source; diamonds mark the six-source mean. Positive values favor the hybrid. The source spread is descriptive, not a confidence interval.}{fig:singlepaired}
\FloatBarrier
Against the primary three-hour control, the hybrid has higher drift, noise and family macro F1 in every one of the 60 cells in each network. Replay is higher in 54 Modena cells and all 60 L-Town cells; pooled F1 is higher in 59 Modena cells and 20 L-Town cells. Hybrid clean FPR is higher in every cell in both networks. Figure~\ref{fig:singlepaired} plots family differences averaged within sources. The frozen summary also retains model and source variation.

Later evidence helps the single classifier too. Its pooled mean increases from 0.8259 to 0.8509 in Modena and from 0.8685 to 0.8977 in L-Town when it is allowed three hours; drift and noise improve in both networks. The compound design gives broader temporal-family coverage under that allowance, particularly in Modena, but does not win on every metric. Since the zero- and three-hour controls are fitted and calibrated separately, the comparison measures the complete configurations at each information allowance.


